
\documentclass{article}
\usepackage{colortbl}
\usepackage{makecell}
\usepackage{multirow}
\usepackage{supertabular}

\begin{document}

\newcounter{utterance}

\centering \large Interaction Transcript for game `ta\_blackjack', experiment `Blackjack{-}v0', episode 5 with BSU{-}SLIM\_\_Qwen3.5{-}9B{-}Base\_\_playpen{-}prm{-}9b{-}best\_of\_n.
\vspace{24pt}

{ \footnotesize  \setcounter{utterance}{1}
\setlength{\tabcolsep}{0pt}
\begin{supertabular}{c@{$\;$}|p{.15\linewidth}@{}p{.15\linewidth}p{.15\linewidth}p{.15\linewidth}p{.15\linewidth}p{.15\linewidth}}
    \# & \multicolumn{2}{c}{Player} && \multicolumn{2}{c}{Game Master} \\
    \hline

    \theutterance \stepcounter{utterance}  
    & & & \multicolumn{4}{p{0.6\linewidth}}{
        \cellcolor[rgb]{0.9,0.9,0.9}{
            \makecell[{{p{\linewidth}}}]{
                \texttt{\tiny{[P1$\langle$GM]}}
                \texttt{You are a competitive game player. Make sure you read the game instructions carefully, and always follow the required format.} \\
\\ 
\texttt{{[}GAME{]} You are playing Blackjack against the dealer.} \\
\texttt{Your goal is to get as close to 21 as possible without going over.} \\
\texttt{On your turn, choose '{[}Hit{]}' to draw another card or '{[}Stand{]}' to hold.} \\
\texttt{J/Q/K = 10 points; A = 11 or 1, whichever is better.} \\
\\ 
\texttt{{[}GAME{]} Hand 1/5} \\
\texttt{Your hand: 3♥, K♦ (Score: 13)} \\
\texttt{Dealer shows: 6♠} \\
            }
        }
    }
     \\ \\

    \theutterance \stepcounter{utterance}  
    & \multicolumn{4}{p{0.6\linewidth}}{
        \cellcolor[rgb]{0.561,0.737,0.561}{
            \makecell[{{p{\linewidth}}}]{
                \texttt{\tiny{[P1$\rangle$GM]}}
                \texttt{**Analysis:**} \\
\texttt{{-} **My Hand:** 3♥ (3 points) + K♦ (10 points) = 13 points.} \\
\texttt{{-} **Dealer's Upcard:** 6♠.} \\
\texttt{{-} **Strategy:** In Blackjack basic strategy, when the dealer shows a 6, they are statistically weaker than showing a 2 through 6 or a 7 through Ace. Hitting on 13 has a decent chance of busting (going over 21) if I draw a 10{-}value card, a Queen, King, or another 10 (roughly 30{-}31\% chance). However, standing on 13 against a dealer's 6 gives the dealer a higher probability of making a hand (since they need to hit until reaching 17, and they are vulnerable to busting themselves).} \\
\texttt{{-} **Decision:** Since my hand is 13 and the dealer has a 6, which is a weak upcard, I should **Hit** to try and improve my total closer to 21 before the dealer plays.} \\
\\ 
\texttt{**{[}Hit{]}**} \\
            }
        }
    }
    & & \\ \\

    \theutterance \stepcounter{utterance}  
    & & & \multicolumn{4}{p{0.6\linewidth}}{
        \cellcolor[rgb]{0.9,0.9,0.9}{
            \makecell[{{p{\linewidth}}}]{
                \texttt{\tiny{[P1$\langle$GM]}}
                \texttt{{[}GAME{]} Hand 1/5} \\
\texttt{Your hand: 3♥, K♦, 6♥ (Score: 19)} \\
\texttt{Dealer shows: 6♠} \\
            }
        }
    }
     \\ \\

    \theutterance \stepcounter{utterance}  
    & \multicolumn{4}{p{0.6\linewidth}}{
        \cellcolor[rgb]{0.561,0.737,0.561}{
            \makecell[{{p{\linewidth}}}]{
                \texttt{\tiny{[P1$\rangle$GM]}}
                \texttt{**Analysis:**} \\
\texttt{*   **My Hand:** 3♥ (3) + K♦ (10) + 6♥ (6) = 19.} \\
\texttt{*   **Dealer's Upcard:** 6♠.} \\
\texttt{*   **Strategy:** I now have 19 points. The dealer shows a 6, which is a weak upcard. In basic strategy, standing on 17 or higher is always the correct play because busting (going over 21) is impossible at this count. Even if the dealer hits their 6 and improves to 17 or 18, I still have a very strong hand at 19.} \\
\\ 
\texttt{**{[}Stand{]}**} \\
            }
        }
    }
    & & \\ \\

    \theutterance \stepcounter{utterance}  
    & & & \multicolumn{4}{p{0.6\linewidth}}{
        \cellcolor[rgb]{0.9,0.9,0.9}{
            \makecell[{{p{\linewidth}}}]{
                \texttt{\tiny{[P1$\langle$GM]}}
                \texttt{{[}GAME{]} Hand 1: you win. Your final 19, Dealer 23.} \\
\texttt{{[}GAME{]} Hand 2/5} \\
\texttt{Your hand: 3♣, 8♥ (Score: 11)} \\
\texttt{Dealer shows: 7♦} \\
            }
        }
    }
     \\ \\

    \theutterance \stepcounter{utterance}  
    & \multicolumn{4}{p{0.6\linewidth}}{
        \cellcolor[rgb]{0.561,0.737,0.561}{
            \makecell[{{p{\linewidth}}}]{
                \texttt{\tiny{[P1$\rangle$GM]}}
                \texttt{**Analysis:**} \\
\texttt{*   **My Hand:** 3♣ (3) + 8♥ (8) = 11 points.} \\
\texttt{*   **Dealer's Upcard:** 7♦.} \\
\texttt{*   **Strategy:** The dealer shows a 7, which is considered a "tough" upcard because they have a relatively high chance of making a standing hand (17{-}20) or hitting into a strong hand. Basic strategy dictates that for hands totaling 11 or 12 against a dealer 7, you should always **Hit**. Standing on 11 carries a significant risk of losing if the dealer makes a 17 or higher, whereas hitting allows you to improve your total with a high probability of ending up with 17, 18, or 19 (which are unbeatable against a dealer 7).} \\
\\ 
\texttt{**{[}Hit{]}**} \\
            }
        }
    }
    & & \\ \\

    \theutterance \stepcounter{utterance}  
    & & & \multicolumn{4}{p{0.6\linewidth}}{
        \cellcolor[rgb]{0.9,0.9,0.9}{
            \makecell[{{p{\linewidth}}}]{
                \texttt{\tiny{[P1$\langle$GM]}}
                \texttt{{[}GAME{]} Hand 2/5} \\
\texttt{Your hand: 3♣, 8♥, Q♥ (Score: 21)} \\
\texttt{Dealer shows: 7♦} \\
            }
        }
    }
     \\ \\

    \theutterance \stepcounter{utterance}  
    & \multicolumn{4}{p{0.6\linewidth}}{
        \cellcolor[rgb]{0.561,0.737,0.561}{
            \makecell[{{p{\linewidth}}}]{
                \texttt{\tiny{[P1$\rangle$GM]}}
                \texttt{**Analysis:**} \\
\texttt{*   **My Hand:** 3♣ (3) + 8♥ (8) + Q♥ (10) = 21 points.} \\
\texttt{*   **Dealer's Upcard:** 7♦.} \\
\texttt{*   **Strategy:** I have reached exactly 21, which is the maximum possible score in Blackjack. Drawing any additional card will result in a bust (score > 21), causing an automatic loss. Since I cannot improve my score, I must stop drawing.} \\
\\ 
\texttt{**{[}Stand{]}**} \\
            }
        }
    }
    & & \\ \\

    \theutterance \stepcounter{utterance}  
    & & & \multicolumn{4}{p{0.6\linewidth}}{
        \cellcolor[rgb]{0.9,0.9,0.9}{
            \makecell[{{p{\linewidth}}}]{
                \texttt{\tiny{[P1$\langle$GM]}}
                \texttt{{[}GAME{]} Hand 2: you win. Your final 21, Dealer 17.} \\
\texttt{{[}GAME{]} Hand 3/5} \\
\texttt{Your hand: 5♥, 9♣ (Score: 14)} \\
\texttt{Dealer shows: J♣} \\
            }
        }
    }
     \\ \\

    \theutterance \stepcounter{utterance}  
    & \multicolumn{4}{p{0.6\linewidth}}{
        \cellcolor[rgb]{0.561,0.737,0.561}{
            \makecell[{{p{\linewidth}}}]{
                \texttt{\tiny{[P1$\rangle$GM]}}
                \texttt{**Analysis:**} \\
\texttt{*   **My Hand:** 5♥ (5) + 9♣ (9) = 14 points.} \\
\texttt{*   **Dealer's Upcard:** J♣ (10 points).} \\
\texttt{*   **Strategy:** The dealer showing a 10{-}value card (J, Q, K, 10) is statistically the most dangerous situation for the player. The dealer has a very high probability of making a strong hand (17, 18, 19, 20, or 21). Standing on 14 against a dealer 10 results in a loss about 55{-}60\% of the time if the dealer doesn't bust. Hitting on 14 gives you a chance to improve to 17, 18, or 19 (which would beat the dealer's likely standing total of 17{-}20) or bust. Since hitting offers the best mathematical expectation to win or push in this specific scenario compared to standing, the correct move is to draw another card.} \\
\\ 
\texttt{**{[}Hit{]}**} \\
            }
        }
    }
    & & \\ \\

    \theutterance \stepcounter{utterance}  
    & & & \multicolumn{4}{p{0.6\linewidth}}{
        \cellcolor[rgb]{0.9,0.9,0.9}{
            \makecell[{{p{\linewidth}}}]{
                \texttt{\tiny{[P1$\langle$GM]}}
                \texttt{{[}GAME{]} Hand 3/5} \\
\texttt{Your hand: 5♥, 9♣, 4♦ (Score: 18)} \\
\texttt{Dealer shows: J♣} \\
            }
        }
    }
     \\ \\

    \theutterance \stepcounter{utterance}  
    & \multicolumn{4}{p{0.6\linewidth}}{
        \cellcolor[rgb]{0.561,0.737,0.561}{
            \makecell[{{p{\linewidth}}}]{
                \texttt{\tiny{[P1$\rangle$GM]}}
                \texttt{**Analysis:**} \\
\texttt{*   **My Hand:** 5♥ (5) + 9♣ (9) + 4♦ (4) = 18 points.} \\
\texttt{*   **Dealer's Upcard:** J♣ (10 points).} \\
\texttt{*   **Strategy:** I have 18 points against a dealer showing 10. In basic strategy, when holding 18 or higher, you should always **Stand**. Even though the dealer has a 10 (a strong upcard), they cannot improve their hand past 20 without busting (since they must hit on 16 or less, and 10+16=26 > 21). My 18 is strong enough to beat the dealer's maximum potential standing score of 20, and I have eliminated the risk of busting by drawing another card.} \\
\\ 
\texttt{**{[}Stand{]}**} \\
            }
        }
    }
    & & \\ \\

    \theutterance \stepcounter{utterance}  
    & & & \multicolumn{4}{p{0.6\linewidth}}{
        \cellcolor[rgb]{0.9,0.9,0.9}{
            \makecell[{{p{\linewidth}}}]{
                \texttt{\tiny{[P1$\langle$GM]}}
                \texttt{{[}GAME{]} Hand 3: you lose. Your final 18, Dealer 20.} \\
\texttt{{[}GAME{]} Hand 4/5} \\
\texttt{Your hand: 7♠, A♦ (Score: 18)} \\
\texttt{Dealer shows: 10♦} \\
            }
        }
    }
     \\ \\

    \theutterance \stepcounter{utterance}  
    & \multicolumn{4}{p{0.6\linewidth}}{
        \cellcolor[rgb]{0.561,0.737,0.561}{
            \makecell[{{p{\linewidth}}}]{
                \texttt{\tiny{[P1$\rangle$GM]}}
                \texttt{**Analysis:**} \\
\texttt{*   **My Hand:** 7♠ (7) + A♦ (11) = 18 points.} \\
\texttt{*   **Dealer's Upcard:** 10♦ (10 points).} \\
\texttt{*   **Strategy:** I have 18 points against a dealer showing 10. Just as in the previous hand of this round, this is a standard situation where standing is the optimal play. The dealer has a 10, which is the strongest possible upcard. While I could hit trying to get to 21, drawing any 10{-}value card (10, J, Q, K) would result in 28 (bust), and drawing a small card (2{-}6) would result in 19 or 20, which are statistically less likely to win against a dealer 10 than simply holding the 18. Basic strategy dictates standing immediately on 18 when facing a dealer 10.} \\
\\ 
\texttt{**{[}Stand{]}**} \\
            }
        }
    }
    & & \\ \\

    \theutterance \stepcounter{utterance}  
    & & & \multicolumn{4}{p{0.6\linewidth}}{
        \cellcolor[rgb]{0.9,0.9,0.9}{
            \makecell[{{p{\linewidth}}}]{
                \texttt{\tiny{[P1$\langle$GM]}}
                \texttt{{[}GAME{]} Hand 4: you win. Your final 18, Dealer 17.} \\
\texttt{{[}GAME{]} Hand 5/5} \\
\texttt{Your hand: J♦, 9♠ (Score: 19)} \\
\texttt{Dealer shows: 7♣} \\
            }
        }
    }
     \\ \\

    \theutterance \stepcounter{utterance}  
    & \multicolumn{4}{p{0.6\linewidth}}{
        \cellcolor[rgb]{0.561,0.737,0.561}{
            \makecell[{{p{\linewidth}}}]{
                \texttt{\tiny{[P1$\rangle$GM]}}
                \texttt{**Analysis:**} \\
\texttt{*   **My Hand:** J♦ (10) + 9♠ (9) = 19 points.} \\
\texttt{*   **Dealer's Upcard:** 7♣.} \\
\texttt{*   **Strategy:** I have 19 points, which is a very strong hand. The dealer shows a 7, which is a "stiff" upcard (meaning the dealer has a good chance of making a 17, 18, 19, or 20). Even if the dealer hits the 7 and improves to a standing 20, my 19 would still result in a push (tie) in most rule sets, or a loss if the dealer hits to an Ace and makes 21. However, drawing another card carries a significant risk of busting (going over 21). Since I already have a score that beats the dealer's initial card and is high enough to compete with their likely final total, there is no benefit to hitting. Standing is the mathematically correct play.} \\
\\ 
\texttt{**{[}Stand{]}**} \\
            }
        }
    }
    & & \\ \\

    \theutterance \stepcounter{utterance}  
    & & & \multicolumn{2}{p{0.3\linewidth}}{
        \cellcolor[rgb]{0.9,0.9,0.9}{
            \makecell[{{p{\linewidth}}}]{
                \texttt{\tiny{[GM$|$GM]}}
                \texttt{role: Player 0} \\
\texttt{invalid\_move: False} \\
\texttt{turn\_count: 8} \\
\texttt{reason: The game has concluded. Final scores: Dealer: 1, You: 4, Draws: 0} \\
\texttt{reward: 0.8} \\
            }
        }
    }
    & & \\ \\

\end{supertabular}
}

\end{document}
